%%%%%%%%%%%%
%
% $Autor: Wings $
% $Datum: 2019-03-05 08:03:15Z $
% $Pfad: SetUp.tex $
% $Version: 4250 $
% !TeX spellcheck = en_GB/de_DE
% !TeX encoding = utf8
% !TeX root = manual 
% !TeX TXS-program:bibliography = txs:///biber
%
%%%%%%%%%%%%

\chapter{Aufbau}
\section{Genereller Aufbau}
Wie in der Skizze \ref{fig:Skizze} zu sehen, besteht der Magic Faucet Fountain aus einem Blumentopf mit montierbaren Deckel. Im Deckel integriert befindet sich ein Acryl-Rohr auf dem der Wasserhahn montiert ist. Eine Pumpe am Boden des Rohres befördert das Wasser nach oben. Der Deckel des Blumentopfes dient zudem als doppelter Boden auf dem die Deko-Steine platziert werden. Im Blumentopf befindet sich zudem ein separates Rohr mit Kunststoff-Schwimmer, dieser dient als Reflektor für den eingebauten Abstandssensor, um den Wasserpegel zu messen. Im Netzteilkasten mit Netzkabel befindet sich die Steuerung, der Hauptschalter und die Stromversorgung des Magic Faucet Fountain.

\section{Wasserführung}
Eine Pumpe befördert das Wasser aus einem Blumentopf durch ein transparentes Acrylrohr nach oben. Das Wasser tritt am Auslass des Hahns aus und strömt nach unten, wobei das Rohr durch den Wasserfluss verdeckt wird.
Der Blumentopf dient als Wasserreservoir und ist mit Ablassbohrungen ausgestattet. Diese ermöglichen den kontrollierten Wasserstand im Behälter und dienen gleichzeitig zur Durchführung der Stromkabel.

\section{Pegelmessung}
Zur Überwachung des Wasserstands ist eine Sensoraufnahme mit Gewinde und Bohrungen im Blumentopfdeckel vorgesehen, in der ein Ultraschallsensor vom Typ HC-SR04 eingebaut ist. Der Blumentopfdeckel trennt zudem die Steine vom Wasser, um eine genauere Messung zu ermöglichen.
Ein Messrohr ist an der Sensoraufnahme befestigt. Innerhalb des Messrohrs befindet sich ein Styropor-Schwimmer, der sich mit steigendem Wasserstand nach oben bewegt. Dadurch kann der Sensor den Wasserstand erfassen. Eine zusätzliche Bohrung im Messrohr dient zur Luftdruckregulierung und für die Kabeldurchführung der Stromleitungen

\section{Elektronik und Steuerung:}
Die Steuerung und Stromversorgung ist in einem separaten Netzteilkasten untergebracht. Dieser enthält:
\begin{enumerate}
	\item Arduino Nano 33 BLE Sense Light zur Verarbeitung der Sensordaten, Bluetoothverbindung und Schaltung des MOSFETs
	\item MOSFET zur Ansteuerung des Laststromkreises der Pumpe
	\item Hauptschalter um das System ein- und auszuschalten
	\item Netzteil mit 2 Ausgangsspannungen, um den Stromkreis für Arduino und Pumpe zu versorgen
	\item Netzkabel zur Stromversorgung
	\item Netz- und Sensorleitung zum Blumentopf
\end{enumerate}
